% GENERATED FILE. Edit canonical lessons, then run npm run generate:books.
% canonical-source-hash: fnv1a64:9db7a2c4a7944bf1
% canonical-lessons: LA-C504-probitas, LA-C504-sincerus, LA-C504-mendax, LA-C504-simulo, LA-C504-dissimulo, LA-R504-review-497-504

\chapter{Honesty and pretence}
\label{ch:la-honesty-and-pretence}
\hlchaptermodality{504}

\begin{chapteropening}
\textbf{By the end of this chapter:} \emph{I can talk about honesty, lies and pretence.}

Complete the last lesson of chapter 504: I can talk about honesty, lies and pretence.
\end{chapteropening}

\section[probitās, probitātis]{probitās, probitātis — honesty, uprightness}

\label{lesson:LA-C504-probitas}



\begin{quote}
Before the new one: say the Latin for charming, then the Latin for playful.
\end{quote}

\subsection*{You'll want to know: probitās, probitātis}
\textbf{probitās, probitātis} — \textquotedblleft{}honesty, uprightness\textquotedblright{}.

\begin{quote}
Omnēs probitātem eius laudant.
\end{quote}

\begin{quote}
\emph{Everyone praises his honesty.}
\end{quote}

Feminine, third declension, from \emph{probus} — upright. English \emph{probity}.

\begin{cousinweb}
\textbf{probitās} is built on \emph{probus} \textquotedblleft{}good, upright\textquotedblright{} (Lewis \& Short), which EtymDB traces to a Proto-Indo-European compound meaning \textquotedblleft{}being in front, excellent\textquotedblright{}. English \emph{probity} comes from it.
\end{cousinweb}

\begin{grammarlens}[title={What you've built}]
One more word for telling, explaining and describing at B1.
\end{grammarlens}

\subsection*{Your turn}
\begin{itemize}
\raggedright
  \item \emph{Say it:} \emph{probitās, probitātis}
  \item \emph{Say it:} \emph{probitās, probitātis}, once more
  \item \emph{Say it:} a sentence of your own with \emph{probitās, probitātis}
  \item \emph{Recall:} say the Latin for charming, then the Latin for playful, then say \emph{probitās, probitātis} again
\end{itemize}

\begin{tcolorbox}[breakable,colback=teal!4,colframe=teal!35!black,title={Before you move on}]
What does probitās, probitātis mean? (\textquotedblleft{}Honesty, uprightness\textquotedblright{}.) Say it once more.
\end{tcolorbox}
\section[sincērus, sincēra, sincērum]{sincērus, sincēra, sincērum — genuine, sincere}

\label{lesson:LA-C504-sincerus}



\begin{quote}
Before the new one: say the Latin for playful, then the Latin for honesty.
\end{quote}

\subsection*{You'll want to know: sincērus, sincēra, sincērum}
\textbf{sincērus, sincēra, sincērum} — \textquotedblleft{}genuine, sincere\textquotedblright{}.

\begin{quote}
Amīcitia nostra sincēra est.
\end{quote}

\begin{quote}
\emph{Our friendship is genuine.}
\end{quote}

Of things, pure and unmixed; of people, honest. English \emph{sincere}.

\begin{cousinweb}
\textbf{sincērus} is a native Latin adjective whose origin is not settled; EtymDB and Lewis \& Short both link its first part with \emph{sim-}, as in \emph{simplex}. English \emph{sincere} comes from it.
\end{cousinweb}

\begin{grammarlens}[title={What you've built}]
One more word for telling, explaining and describing at B1.
\end{grammarlens}

\subsection*{Your turn}
\begin{itemize}
\raggedright
  \item \emph{Say it:} \emph{sincērus, sincēra, sincērum}
  \item \emph{Say it:} \emph{sincērus, sincēra, sincērum}, once more
  \item \emph{Say it:} a sentence of your own with \emph{sincērus, sincēra, sincērum}
  \item \emph{Recall:} say the Latin for playful, then the Latin for honesty, then say \emph{sincērus, sincēra, sincērum} again
\end{itemize}

\begin{tcolorbox}[breakable,colback=teal!4,colframe=teal!35!black,title={Before you move on}]
What does sincērus, sincēra, sincērum mean? (\textquotedblleft{}Genuine, sincere\textquotedblright{}.) Say it once more.
\end{tcolorbox}
\section[mendāx, mendācis]{mendāx, mendācis — lying; a liar}

\label{lesson:LA-C504-mendax}



\begin{quote}
Before the new one: say the Latin for honesty, then the Latin for genuine.
\end{quote}

\subsection*{You'll want to know: mendāx, mendācis}
\textbf{mendāx, mendācis} — \textquotedblleft{}lying; a liar\textquotedblright{}.

\begin{quote}
Nēmō mendācī crēdit.
\end{quote}

\begin{quote}
\emph{No one believes a liar.}
\end{quote}

One ending: \emph{mendāx} for all genders. As a noun, a liar. \emph{Mendācium}, already known, is the lie.

\begin{cousinweb}
\textbf{mendāx} belongs with \emph{mentior} \textquotedblleft{}to lie\textquotedblright{} and \emph{mendācium} \textquotedblleft{}a lie\textquotedblright{} (Lewis \& Short). English \emph{mendacious} comes from it.
\end{cousinweb}

\begin{grammarlens}[title={What you've built}]
One more word for telling, explaining and describing at B1.
\end{grammarlens}

\subsection*{Your turn}
\begin{itemize}
\raggedright
  \item \emph{Say it:} \emph{mendāx, mendācis}
  \item \emph{Say it:} \emph{mendāx, mendācis}, once more
  \item \emph{Say it:} a sentence of your own with \emph{mendāx, mendācis}
  \item \emph{Recall:} say the Latin for honesty, then the Latin for genuine, then say \emph{mendāx, mendācis} again
\end{itemize}

\begin{tcolorbox}[breakable,colback=teal!4,colframe=teal!35!black,title={Before you move on}]
What does mendāx, mendācis mean? (\textquotedblleft{}Lying; a liar\textquotedblright{}.) Say it once more.
\end{tcolorbox}
\section[simulō, simulāre]{simulō, simulāre — to pretend, imitate}

\label{lesson:LA-C504-simulo}



\begin{quote}
Before the new one: say the Latin for genuine, then the Latin for lying.
\end{quote}

\subsection*{You'll want to know: simulō, simulāre}
\textbf{simulō, simulāre} — \textquotedblleft{}to pretend, imitate\textquotedblright{}.

\begin{quote}
Servus morbum simulat, quod labōrāre nōn vult.
\end{quote}

\begin{quote}
\emph{The slave pretends to be ill, because he does not want to work.}
\end{quote}

\emph{simulō, simulāre, simulāvī, simulātum}: to pretend that something is so. English \emph{simulate}.

\begin{cousinweb}
\textbf{simulō} is built on \emph{similis} \textquotedblleft{}like\textquotedblright{} (Lewis \& Short; EtymDB): to make a thing look like another. English \emph{simulate} comes from it.
\end{cousinweb}

\begin{grammarlens}[title={What you've built}]
One more word for telling, explaining and describing at B1.
\end{grammarlens}

\subsection*{Your turn}
\begin{itemize}
\raggedright
  \item \emph{Say it:} \emph{simulō, simulāre}
  \item \emph{Say it:} \emph{simulō, simulāre}, once more
  \item \emph{Say it:} a sentence of your own with \emph{simulō, simulāre}
  \item \emph{Recall:} say the Latin for genuine, then the Latin for lying, then say \emph{simulō, simulāre} again
\end{itemize}

\begin{tcolorbox}[breakable,colback=teal!4,colframe=teal!35!black,title={Before you move on}]
What does simulō, simulāre mean? (\textquotedblleft{}To pretend, imitate\textquotedblright{}.) Say it once more.
\end{tcolorbox}
\section[dissimulō, dissimulāre]{dissimulō, dissimulāre — to hide, disguise (a feeling)}

\label{lesson:LA-C504-dissimulo}



\begin{quote}
Before the new one: say the Latin for lying, then the Latin for to pretend.
\end{quote}

\subsection*{You'll want to know: dissimulō, dissimulāre}
\textbf{dissimulō, dissimulāre} — \textquotedblleft{}to hide, disguise (a feeling)\textquotedblright{}.

\begin{quote}
Puella dolōrem suum dissimulat et rīdet.
\end{quote}

\begin{quote}
\emph{The girl hides her pain and smiles.}
\end{quote}

\emph{dissimulō, dissimulāre, dissimulāvī, dissimulātum}: to pretend that something is not so. English \emph{dissimulate}.

\begin{cousinweb}
\textbf{dissimulō} is built on \emph{dissimilis} \textquotedblleft{}unlike\textquotedblright{} (Lewis \& Short; EtymDB). English \emph{dissimulate} comes from it, and \emph{dissemble} is a reshaped form.
\end{cousinweb}

\begin{grammarlens}[title={What you've built}]
One more word for telling, explaining and describing at B1.
\end{grammarlens}

\subsection*{Your turn}
\begin{itemize}
\raggedright
  \item \emph{Say it:} \emph{dissimulō, dissimulāre}
  \item \emph{Say it:} \emph{dissimulō, dissimulāre}, once more
  \item \emph{Say it:} a sentence of your own with \emph{dissimulō, dissimulāre}
  \item \emph{Recall:} say the Latin for lying, then the Latin for to pretend, then say \emph{dissimulō, dissimulāre} again
\end{itemize}

\begin{tcolorbox}[breakable,colback=teal!4,colframe=teal!35!black,title={Before you move on}]
What does dissimulō, dissimulāre mean? (\textquotedblleft{}To hide, disguise (a feeling)\textquotedblright{}.) Say it once more.
\end{tcolorbox}
\section[(dialogue)]{Words from chapters 497–504 — a review}

\label{lesson:LA-R504-review-497-504}



\begin{quote}
Say the Latin for anxiety and to distress.
\end{quote}

\subsection*{Your turn}
Say these aloud:

\begin{itemize}
\raggedright
  \item anxiety, to distress, distress of mind, annoyance, to torment
  \item to despair, despair, to trust, to distrust, trusting
  \item to delight, delight, pleasure, to amuse, to please
  \item desire, eager, greedy, to envy, a rival
  \item to move deeply, to affect, a mood, calm, free from care
  \item a joke, to joke, humorous, funny, witty
  \item nature, natural disposition, lively, charming, playful
  \item honesty, genuine, lying, to pretend, to hide
\end{itemize}

\begin{tcolorbox}[breakable,colback=teal!4,colframe=teal!35!black,title={Before you move on}]
Anxiety? (\textbf{sollicitūdō, sollicitūdinis}.) To distress? (\textbf{angō, angere}.) Distress of mind? (\textbf{aegritūdō, aegritūdinis}.) Annoyance? (\textbf{molestia, molestiae}.) To torment? (\textbf{cruciō, cruciāre}.) To despair? (\textbf{dēspērō, dēspērāre}.) Despair? (\textbf{dēspērātiō, dēspērātiōnis}.) To trust? (\textbf{cōnfīdō, cōnfīdere}.) To distrust? (\textbf{diffīdō, diffīdere}.) Trusting? (\textbf{crēdulus, crēdula, crēdulum}.) To delight? (\textbf{dēlectō, dēlectāre}.) Delight? (\textbf{dēlectātiō, dēlectātiōnis}.) Pleasure? (\textbf{voluptās, voluptātis}.) To amuse? (\textbf{oblectō, oblectāre}.) To please? (\textbf{placeō, placēre}.) Desire? (\textbf{cupīdō, cupīdinis}.) Eager? (\textbf{cupidus, cupida, cupidum}.) Greedy? (\textbf{avidus, avida, avidum}.) To envy? (\textbf{invideō, invidēre}.) A rival? (\textbf{rīvālis, rīvālis}.) To move deeply? (\textbf{commoveō, commovēre}.) To affect? (\textbf{afficiō, afficere}.) A mood? (\textbf{affectus, affectūs}.) Calm? (\textbf{tranquillus, tranquilla, tranquillum}.) Free from care? (\textbf{sēcūrus, sēcūra, sēcūrum}.) A joke? (\textbf{iocus, iocī}.) To joke? (\textbf{iocor, iocārī}.) Humorous? (\textbf{iocōsus, iocōsa, iocōsum}.) Funny? (\textbf{rīdiculus, rīdicula, rīdiculum}.) Witty? (\textbf{facētus, facēta, facētum}.) Nature? (\textbf{nātūra, nātūrae}.) Natural disposition? (\textbf{indolēs, indolis}.) Lively? (\textbf{alacer, alacris, alacre}.) Charming? (\textbf{lepidus, lepida, lepidum}.) Playful? (\textbf{lascīvus, lascīva, lascīvum}.) Honesty? (\textbf{probitās, probitātis}.) Genuine? (\textbf{sincērus, sincēra, sincērum}.) Lying? (\textbf{mendāx, mendācis}.) To pretend? (\textbf{simulō, simulāre}.) To hide? (\textbf{dissimulō, dissimulāre}.)
\end{tcolorbox}
